% !TEX root = Main.tex
\section{Lower bound in the general adversarial setting} %(Theorem~\ref{th:LOWadPRO}) }
\label{app:prooflowadPRO}
%
	\setcounter{scratchcounter}{\value{theorem}}\setcounter{theorem}{\the\numexpr\getrefnumber{th:LOWadPRO}-1}
\toto*
\setcounter{theorem}{\the\numexpr\value{scratchcounter}}
\begin{proof}%[Theorem~\ref{th:LOWadPRO}]
Fix the reference gaps
$0<\gap_1=\gap_2\leq\cdots\leq\gap_{\nArms}\leq1/8$,
a learner, and a positive integer horizon $\timeHorizon$ satisfying
Theorem~\ref{th:LOWadPRO}. Write
$H=\complexityUnif=\nArms/\gap_1^2$ for the reference complexity.
The learner may know the horizon. The constructions and selected arm
below may depend on the learner and horizon, but not on its realized
random choices. Set
\[
m=\lfloor\timeHorizon/2\rfloor,\qquad \beta=m/\timeHorizon.
\]
The horizon assumption gives
\[
\timeHorizon\gap_1^2\geq128\log(128\nArms)\geq1.
\]
Since $\gap_1\leq1/8$, we have $\timeHorizon\geq64$, and therefore
$1\leq m<\timeHorizon$. Moreover,
\[
\frac{11}{24}\leq\frac12-\frac{1}{2\timeHorizon}
\leq\beta\leq\frac12,
\]
where the first inequality uses $\timeHorizon\geq12$.

Consider a base problem \basePro{} with independent Bernoulli rewards
of means $\mean_1^{\basePro}=1/2$ and
$\mean_j^{\basePro}=1/2-\gap_j$ for $j\in\setArmsmo$.
Let $\pullsNumber_j(m)$ be the number of pulls of arm $j$ through
round $m$. Since the expected pull counts of arms in $\setArmsmo$
sum to at most $m$, choose a fixed arm $\refarm\in\setArmsmo$ such that
\begin{equation}\label{eq:notSoMuchPullsAD}
\Exp_{\basePro}\left[\pullsNumber_{\refarm}(m)\right]
\leq\frac{m}{\nArms-1}\leq\frac{\timeHorizon}{\nArms}.
\end{equation}
The last inequality follows from $m\leq\timeHorizon/2$ and
$\nArms\leq2(\nArms-1)$ for $\nArms\geq2$.
This arm is chosen from expectations under \basePro{}, before either
adversarial table is sampled. The initial phase is the prefix
$\earlyPhase=\{1,\ldots,m\}$.

Set
\[
c=\gap_{\refarm}-\gap_1.
\]
Since $\gap_1\leq\gap_{\refarm}\leq1/8$, we have $0\leq c\leq1/8$.
Construct two distributions $\advPro_1$ and $\advPro_2$ over complete
reward tables, sampled before interaction and independently of the
learner's internal randomness. In $\advPro_1$, the prefix follows
\basePro{} on all arms. In $\advPro_2$, the prefix also follows
\basePro{}, except that arm $\refarm$ has Bernoulli mean
$1/2-\gap_{\refarm}+2\gap_1$. All prefix entries are independent
across arms and rounds. In both environments the second phase has
$m$ rounds and is identical and deterministic: for $m<t\leq2m$,
\[
\gainVector_{\refarm,t}=c,\qquad
\gainVector_{j,t}=0\quad(j\neq\refarm).
\]
For $2m<t\leq\timeHorizon$, set $\gainVector_{j,t}=0$ for every arm
$j$. Thus even horizons give the original two equal phases, while
odd horizons append one all-zero final round. These rewards lie in
$[0,1]$. The original learner, configured for horizon $\timeHorizon$,
chooses an arm in every round, including the optional last round,
and recommends an arm only after round $\timeHorizon$.
Let $\Pro_\ell$ and $\Exp_\ell$ denote the full table-and-learner
probabilities and expectations in $\advPro_\ell$, for $\ell\in\{1,2\}$.
On tables with tied maximizing arms, define $\bestArm_{\gainVector}$
as the smallest maximizing index.

Write $M_j^\ell=\Exp_\ell[\cumulGain_j]$, where
$\cumulGain_j=\sum_{t=1}^{\timeHorizon}\gainVector_{j,t}$.
The identity $mc=m(\gap_{\refarm}-\gap_1)$ gives the exact expected
totals below; the optional final zero round contributes nothing:
\begin{align*}
M_1^1=M_1^2&=m/2,\\
M_j^1=M_j^2&=m(1/2-\gap_j)
\quad(j\notin\{1,\refarm\}),\\
M_{\refarm}^1
&=m(1/2-\gap_{\refarm})+mc
=m(1/2-\gap_1),\\
M_{\refarm}^2
&=m(1/2-\gap_{\refarm}+2\gap_1)+mc
=m(1/2+\gap_1).
\end{align*}
The common deterministic second phase places arm $\refarm$ just below
arm $1$ in the first expected table and just above it in the second. The expected average
gap between them is $\beta\gap_1$, irrespective of the original
size of $\gap_{\refarm}$, while the environments differ only in the
initial rewards of the lightly sampled arm.

For $\ell\in\{1,2\}$ define the full-table good event
\[
\eventc_\ell=\left\{\forall j\in\setArms,\quad
|\cumulGain_j-M_j^\ell|\leq\frac{\timeHorizon\gap_1}{16}\right\}.
\]
Only the first $m$ rewards of each arm are random. Two-sided Hoeffding
and a union bound over the complete arm sequences give
\begin{equation}\label{eq:goodEventGenAd}
\Pro_\ell(\eventc_\ell^c)
\leq2\nArms\exp\left(-\frac{2(\timeHorizon\gap_1/16)^2}{m}\right)
\leq2\nArms\exp\left(-\frac{\timeHorizon\gap_1^2}{64}\right),
\end{equation}
since $m\leq\timeHorizon/2$. These events involve the full reward
table and are generally not determined by the learner's observations.

On either good event, every normalized pairwise difference deviates
from its expectation by at most $\gap_1/8$.
On $\eventc_1$, the expected differences from arm $1$ are
$\beta\gap_1$ for arm $\refarm$ and $\beta\gap_j$ for
$j\notin\{1,\refarm\}$. Since
$\beta-1/8\geq1/3$, $\beta+1/8\leq5/8$, and $\gap_1\leq\gap_j$,
\[
\frac{\gap_1}{3}\leq
\frac{\cumulGain_1-\cumulGain_{\refarm}}{\timeHorizon}
\leq\frac{5\gap_1}{8},\qquad
\frac{\gap_j}{3}\leq
\frac{\cumulGain_1-\cumulGain_j}{\timeHorizon}
\leq\frac{5\gap_j}{8}
\quad(j\notin\{1,\refarm\}).
\]
All these differences are positive, so arm $1$ is uniquely best.
Its own gap is the minimum of these differences; the lower bounds
and the upper bound from arm $\refarm$ give
\[
\frac{\gap_1}{3}\leq\gap_1^{\gainVector}\leq\frac{5\gap_1}{8}.
\]
Thus every label satisfies the defining bounds for $\Gaps_3$, with
$\refarm$ as the exceptional label, and every table in $\eventc_1$
belongs to $\Gaps_3$.

On $\eventc_2$, the expected differences from arm $\refarm$ are
$\beta\gap_1$ for arm $1$ and $\beta(\gap_j+\gap_1)$ for
$j\notin\{1,\refarm\}$. Consequently,
\[
\frac{\gap_1}{3}\leq
\frac{\cumulGain_{\refarm}-\cumulGain_1}{\timeHorizon}
\leq\frac{5\gap_1}{8},\qquad
\frac{\gap_j}{3}\leq
\frac{\cumulGain_{\refarm}-\cumulGain_j}{\timeHorizon}
\leq\frac{9\gap_j}{8}
\quad(j\notin\{1,\refarm\}).
\]
For the latter lower bound,
$\beta(\gap_j+\gap_1)-\gap_1/8
\geq(\beta-1/8)\gap_j\geq\gap_j/3$.
For the upper bound, $\beta\leq1/2$ and $\gap_1\leq\gap_j$ give
$\beta(\gap_j+\gap_1)+\gap_1/8\leq9\gap_j/8$.
Arm $\refarm$ is uniquely best, and its own gap, the minimum of
its differences from every other arm, satisfies
\[
\frac{\gap_1}{3}\leq\gap_{\refarm}^{\gainVector}\leq\frac{5\gap_1}{8}.
\]
The upper bound uses the difference from arm $1$. Label $1$ meets
its ordinary reference-gap requirement, all other labels except
$\refarm$ meet theirs, and $\refarm$ meets the exceptional requirement.
Every table in $\eventc_2$ therefore also belongs to $\Gaps_3$.

Let $P_1,P_2$ be the observed-transcript laws from
Lemma~\ref{l:transcriptChangeMeas} under $\advPro_1,\advPro_2$,
distinguished from the full laws $\Pro_1,\Pro_2$.
Apply the lemma with $A=\advPro_1$, $B=\advPro_2$,
$\mathcal T=\{1,\ldots,m\}$, and
\[
u=\frac12-\gap_{\refarm},\qquad
v=\frac12-\gap_{\refarm}+2\gap_1.
\]
All table entries are independent of one another and of the learner's
randomness. The common distributions have finite support, including
the identical deterministic second phase and optional last round,
and $u,v\in[3/8,5/8]\subset(0,1)$. This entire continuation contributes
no additional likelihood-ratio factors.
The same learner is run with the same horizon.
The lemma gives
\[
D(P_1\Vert P_2)
=\Exp_1[\pullsNumber_{\refarm}(m)]\operatorname{kl}(u,v)
=\Exp_{\basePro}[\pullsNumber_{\refarm}(m)]\operatorname{kl}(u,v).
\]
The second equality holds because $\advPro_1$ and \basePro{} have
identical interaction laws through the prefix, including the learner's
adaptive actions. No law is conditioned on a good-table event.

Using $\log z\leq z-1$,
\[
\operatorname{kl}(u,v)
\leq u\left(\frac{u}{v}-1\right)
+(1-u)\left(\frac{1-u}{1-v}-1\right)
=\frac{(u-v)^2}{v(1-v)}.
\]
Here $|u-v|=2\gap_1$ and $v(1-v)\geq15/64$, so
$\operatorname{kl}(u,v)\leq(256/15)\gap_1^2$.
Set $x=\timeHorizon/H=\timeHorizon\gap_1^2/\nArms>0$.
Equation~\ref{eq:notSoMuchPullsAD} yields
\begin{equation}\label{eq:transcriptKLGenAd}
D(P_1\Vert P_2)
\leq\frac{256}{15}\gap_1^2\frac{\timeHorizon}{\nArms}
=\frac{256}{15}x\leq32x.
\end{equation}
For the transcript event $W=\{\recomArm_{\timeHorizon}=1\}$,
the testing bound in Equation~\ref{eq:transcriptTesting} now gives
\[
P_1(\recomArm_{\timeHorizon}\neq1)+P_2(\recomArm_{\timeHorizon}=1)
\geq\frac12\exp(-D(P_1\Vert P_2))\geq\frac12\exp(-32x).
\]
The learner may recommend any arm; recommending arm $1$ is simply
one error event in the second environment on its good tables.

Write $\varepsilon=\exp(-32x)$ and
$b_\ell=\Pro_\ell(\eventc_\ell^c)$.
Both theorem assumptions are needed to compare the bad-table
probabilities with this testing bound. Equation~\ref{eq:goodEventGenAd}
and $\nArms\exp(-\timeHorizon\gap_1^2/128)\leq1/128$ give
\begin{align*}
b_\ell
&\leq2\nArms\exp(-\timeHorizon\gap_1^2/64)\\
&=2\left[\nArms\exp(-\timeHorizon\gap_1^2/128)\right]
\exp(-\timeHorizon\gap_1^2/128)\\
&\leq\frac1{64}\exp(-\timeHorizon\gap_1^2/128)
\leq\frac{\varepsilon}{64}.
\end{align*}
The last inequality uses $\nArms\geq4096$, which implies
$32/\nArms\leq1/128$. Thus bad tables cost only a fixed fraction
of the exponentially small testing bound.

For a fixed table $\gainVector$, let
$\ProErr_{\advPro(\gainVector)}(\timeHorizon)$ denote error over the
learner's internal randomness, using the best-arm convention above.
Define its error mass on good tables by
\[
R_\ell=\Exp_{\gainVector\sim\advPro_\ell}
\left[\ProErr_{\advPro(\gainVector)}(\timeHorizon)
\indicator{\gainVector\in\eventc_\ell}\right].
\]
The unique best arms on the good events give
\begin{align*}
R_1&=\Pro_1(\{\recomArm_{\timeHorizon}\neq1\}\cap\eventc_1),\\
R_2&=\Pro_2(\{\recomArm_{\timeHorizon}\neq\refarm\}\cap\eventc_2)
\geq\Pro_2(\{\recomArm_{\timeHorizon}=1\}\cap\eventc_2).
\end{align*}
Subtracting the bad-table probabilities therefore yields
\begin{equation}\label{eq:goodErrorGenAd}
\begin{aligned}
R_1+R_2
&\geq P_1(\recomArm_{\timeHorizon}\neq1)
+P_2(\recomArm_{\timeHorizon}=1)-b_1-b_2\\
&\geq\frac{\varepsilon}{2}-\frac{2\varepsilon}{64}
=\frac{15\varepsilon}{32}.
\end{aligned}
\end{equation}
All joint events involving $\eventc_\ell$ use the full law $\Pro_\ell$;
no independence from the recommendation is required.

Each table distribution has finite support: only its $m\nArms$
prefix entries are random Bernoulli variables and its suffix is fixed.
Each good subset of that support is nonempty, since
$b_\ell\leq\varepsilon/64<1$. Choose a deterministic table
$\gainVector^\ell$ maximizing
$\ProErr_{\advPro(\gainVector)}(\timeHorizon)$ over the good support of
$\advPro_\ell$. Then
\[
R_\ell\leq\Pro_\ell(\eventc_\ell)
\ProErr_{\advPro(\gainVector^\ell)}(\timeHorizon)
\leq\ProErr_{\advPro(\gainVector^\ell)}(\timeHorizon).
\]
Equation~\ref{eq:goodErrorGenAd} implies
\[
\max_{\ell\in\{1,2\}}
\ProErr_{\advPro(\gainVector^\ell)}(\timeHorizon)
\geq\frac{R_1+R_2}{2}
\geq\frac{15}{64}\exp\left(-\frac{32\timeHorizon}{H}\right).
\]
Both tables belong to $\Gaps_3$ and have different unique best arms,
$1$ and $\refarm$. Since $15/64\geq1/128$, this bound implies
\[
\max_{\ell\in\{1,2\}}\ProErr_{\gainVector^\ell}(\timeHorizon)
\geq\min\left\{\frac1{128}\exp\left(-\frac{32\timeHorizon}
{\complexityUnif}\right),\frac1{32}\right\},
\]
which is the claimed conclusion. The extraction uses finite-support
maximization and requires neither conditional probabilities on null
events nor a coupling of the two table draws.
\end{proof}%
